\chapter{Coarse spaces}
\label{chap:coarse-spaces}

The project contains a copy of S.~Ghobash's coarse-space library (see \texttt{ATTRIBUTION.md}).
It proves a \v{S}varc--Milnor statement as a coarse equivalence, without constants, word metric
or finite generation. Nothing yet connects it to the metric notions above.

\inputleannode{def:coarse-space}
\inputleannode{def:coarsely-close}
\inputleannode{def:coarse-map}
\inputleannode{def:coarse-equiv}
\inputleannode{def:metric-coarse-space}
\inputleannode{lem:coarsely-proper-smul}
\inputleannode{thm:coarse-svarc-milnor}

\section*{Candidate bridges}

These have no consumer yet; they become ready when a node uses them.

\begin{lemma}
  \label{lem:finite-dist-iff-coarsely-close}
  \notready
  \uses{def:finite-dist, def:coarsely-close, def:metric-coarse-space}
  Two maps into a pseudometric space have finite distance if and only if they are coarsely
  close for the metric coarse structure.
  \plannedlean{Metric.hasFiniteDist_iff_isCoarselyClose}
\end{lemma}
\begin{proof}
  Unfold the controlled sets of the metric coarse structure.
\end{proof}

\begin{lemma}
  \label{lem:qi-coarse-equiv}
  \notready
  \uses{def:qi, def:coarse-equiv, def:metric-coarse-space, lem:finite-dist-iff-coarsely-close}
  A quasi-isometry together with a quasi-inverse is a coarse equivalence for the metric coarse
  structures.
  \plannedlean{Metric.IsQuasiIsometry.toCoarseEquiv}
\end{lemma}
\begin{proof}
  Coarse Lipschitz maps are controlled, coarse antilipschitz maps are coarsely proper, and maps
  at finite distance are coarsely close.
\end{proof}
